\section{Security Analysis}
\label{sec:secannalysis}

\subsection{P1 Strong Microarchitectural Isolation}
MI6 already enforces P1 through partitioning and flushing 
(See Section \ref{sec:hardwaremechanisms}). 
Therefore, we only need to update the security analysis to cover the new mechanisms we introduced.

We detail strong isolation arguments between each security domains,
namely the enclave, the enclave mini-SM, the OS (untrusted software) and the main SM. We also consider the boundary between the OS and the bootloader.

Note that we are not susceptible to Meltdown-style attacks, as the processor does not speculate across privilege levels.
For instance, we always perform security checks even on speculative memory accesses and do not forward illegal results in the pipeline.
Because we do not consider multithreading across security domains, any state above the L1 cache cannot be shared concurrently between two different security domains.
Finally, our processor does not include a prefetcher so we are not susceptible to Spectre variants exploiting them.

\textbf{LLC Dynamic Partitioning}
Our scheme enforces strong isolation between LLC partitions by guaranteeing that 1) the set-index ranges are disjoint and 2) LLC slices are properly sanitized when reshaped.
The SM enforces the first property by ensuring that the sum of all LLC slices fits in the LLC and by enforcing that the base index of each range is equal to the sum of previous range sizes.
The second property is enforced by our double-buffering policy that flushes the LLC before enforcing a new partitioning policy (see Section \ref{sec:dynamiccache}).

\textbf{Enclave/OS Boundary}
Ignoring shared memory for a minute, our enclave benefits from the strong isolation property provided by MI6.
Thanks to physical memory isolation and the partitioning of the memory hierarchy, the enclave and the OS do not share any microarchitectural state.
Because memory and the LLC are still partitioned across context switches, their state is never shared.
When freeing a memory region and the associated cache partition, the SM overwrites the entire region with 0's, effectively purging the LLC partition simultaneously.
For the L1 and the core state, the SM purges every shared structure on context switch.

Nevertheless, our secure shared memory mechanism introduces shared state between the two security domains. 
The important point here is to remember that changes in the newly-shared microarchitecture cannot be performed speculatively:
microarchitectural changes are tied to an observable architectural state change.
Observations possible through the shared microarchitecture are thus strict equivalents to the observations possible through the accesses to shared memory.
Because in our threat model we consider these accesses to be public, no extra information can be leaked through the shared microarchitecture. 
Secrets can still be leaked through the timing of these public interactions, but as we explain in Section~\ref{sec:threadmodel} such attacks could be mounted even after placing the attacker and the victim on two different machines connected through a network.
As a result, we effectively mitigate any transient execution attacks that exploit microarchitectural side channels to leak secrets, which include all Spectre variants besides NETSpectre \cite{schwarz2019netspectre} and achieve our strong microarchitectural isolation property.

\textbf{Enclave/mini-SM Boundary}
In Citadel, we relax the conservative flushing policy of MI6.
We consider the mini-SM as an extension of the security domain of the enclave.
This means we will need to take extra care when interacting with other security domains, but for simple system calls that can be directly served by the mini-SM, we do not need to perform any flushing of the shared microarchitecture.
We evaluate the performance gain of this new policy in Section \ref{sec:eval}.

\textbf{mini-SM/Enclave Boundary}
Here again, we consider the mini-SM as an extension of the security domain of the enclave.
First, like in MI6, speculation is deactivated in the mini-SM, which means that even if it has access to the entire physical address space, it will not access any data speculatively.
Second, we observe that the mini-SM does not perform any sensitive operations (such as keyed cryptography) and does not access any piece of information that should be kept secret from the enclave.
As a result, we do not worry about side channels leaking information from the mini-SM to the enclave as no secret information is accessed.
In particular, we do not need to flush shared state (L1, etc.) on these specific context switches.

\textbf{mini-SM/OS Boundary}
Like in MI6, the mini-SM is hosted within the enclave memory region and its associated LLC partition. 
Its code is trusted, run non-speculatively, and only performs explicit memory accesses to that region with the exception of accesses to a small shared SM state. 
These accesses are public and cannot be used to mount microarchitectural side channels.
Because the mini-SM now shares microarchitectural sate with the enclave, we do not relax the flushing policy when context switching from the mini-SM to the OS and we benefit from MI6 strong isolation.

\textbf{SM/OS Bootloader/OS Boundary}
Similarly to the mini-SM, the SM and the bootloader are trusted and do not speculate on memory accesses but unlike the mini-SM, they might touch sensitive information such as cryptographic attestation keys or enclave secrets.
Therefore, like in MI6, the SM and the bootloader use the same dedicated private memory region and use their own private section of the LLC.
For the bootloader, we flush all microarchitectural state after initialization.
Nonetheless, we can relax the MI6 flushing policy on context switches after SM calls that do not touch sensitive information as no secret can be accessed covertly.
Moreover, instead of flushing all data structures, we also make it possible for the SM to neither use nor train a specific structure to save the OS from a cold start.
In this case, strong isolation is still enforced as no microarchitectural change is performed before re-entering the OS security domain.

Finally, note that in the case of the SM and the bootloader, even if this consideration is outside of our general threat model, we can also make sure their sensitive operations are constant time.
Because their code is executed non-speculatively, and these sensitive functions are written in a constant-time fashion and microarchitectural state is flushed before their execution, the SM and bootloader are protected against timing attacks such as NETSpectre \cite{schwarz2019netspectre} that use attacker-provided inputs to carry out transient execution attacks and leak secrets through timing of public interaction (like the completion time of a system call).

\textbf{OS/Enclave, OS/SM-mini and OS/SM Boundary}
We do not protect the OS against microarchitectural side channels.

\subsection{Other security properties}

\textbf{I2 Runtime Integrity}
Enclave integrity is guaranteed by our physical memory partitioning scheme that ensures that no adversary can modify any code or data for the protected security domains (security monitor and enclaves).

\textbf{A1 OS Availability}
We make sure that at any time, the OS can send an inter-process interrupt to an enclave and reclaim its resources.
Finally, a user-level application cannot directly access the security monitor, but can only interact with it through the SM kernel module.
As a result, the security monitor is always available to the OS.

\textbf{I1 Authenticity}
Enclave authenticity is achieved by our remote attestation protocol.
We use a simpler variant of the remote attestation protocol described in Sanctum \cite{costan2016sanctum} whose security is analyzed in \cite{lebedev2018secure}.